\chapter{小模型与端侧部署}
性能采样接口参考PyTorch官方剖析文档\cite{profiler}，采样开销与普通推理时延分开记录。


端侧工作先证明软件栈与量化存储可用，再做同输入精度对照，最后测时延和联合服务。恢复后的自有板子与10月借用环境不混为一套版本，速度和准确率也不混为一个结论。

\section{Jetson恢复与独立GPU环境}

\label{stage:s21}

恢复板端系统，创建独立推理环境，按JetPack与CUDA版本安装适配的PyTorch、量化库和低秩适配工具，并逐步验证GPU计算与4位加载。

\subsection{结论}\label{s21-detail-1}

JetPack6.2.1/CUDA12.6的Jetson专用PyTorch、torchvision与bitsandbytes已在ali-jetson安装，实际CUDA计算、NF4矩阵乘法、预量化存储跨机器加载、PEFT直接加载均通过。无须源码编译bitsandbytes。下一阶段先2B，4B保留在后。

\FloatBarrier\begingroup\small\setlength{\tabcolsep}{2.5pt}
\begin{longtable}[]{@{}
  >{\raggedright\arraybackslash}p{(\linewidth - 4\tabcolsep) * \real{0.2449}}
  >{\raggedright\arraybackslash}p{(\linewidth - 4\tabcolsep) * \real{0.3094}}
  >{\raggedright\arraybackslash}p{(\linewidth - 4\tabcolsep) * \real{0.4456}}@{}}
\caption{结论：设置与结果}\label{tab:s21-1}\tabularnewline
\toprule\noalign{}
\begin{minipage}[b]{\linewidth}\raggedright
组件
\end{minipage} & \begin{minipage}[b]{\linewidth}\raggedright
版本/来源
\end{minipage} & \begin{minipage}[b]{\linewidth}\raggedright
检查
\end{minipage} \\
\midrule\noalign{}
\endfirsthead
\toprule\noalign{}
\begin{minipage}[b]{\linewidth}\raggedright
组件
\end{minipage} & \begin{minipage}[b]{\linewidth}\raggedright
版本/来源
\end{minipage} & \begin{minipage}[b]{\linewidth}\raggedright
检查
\end{minipage} \\
\midrule\noalign{}
\endhead
\bottomrule\noalign{}
\endlastfoot
torch & 2.8.0，Jetson AI Lab jp6/cu126 & CUDA12.6，设备Orin，SM8.7，矩阵乘法通过 \\
torchvision & 0.23.0，同上 & CUDA NMS通过 \\
bitsandbytes & 0.48.0.dev0+ff389db，同上 & NF4双重量化、4bit Linear前向通过 \\
transformers & 4.57.6 & Qwen3-VL模型类和AutoProcessor可导入，未加载权重 \\
PEFT & 0.18.1 & 在预量化小模型上直接加载适配器通过，未合并 \\
accelerate & 1.12.0 & 依赖安装完成，pip check通过 \\
\end{longtable}\endgroup

torch/vision/bnb均来自https://pypi.jetson-ai-lab.io/jp6/cu126，而非通用aarch64 torch。安装报告记录下载URL与SHA256；后续HF依赖通过约束固定这三个包及NumPy1.26.4，避免解析器替换Jetson构建。系统Python未修改。

\subsection{内存与运行状态}\label{s21-detail-2}

\begin{itemize}
\tightlist
\item
  已从graphical.target切换为multi-user.target并设为默认；恢复桌面可用sudo systemctl set-default graphical.target、sudo systemctl isolate graphical.target。
\item
  板上无NVMe。用户随后明确允许SD卡swap，已在/swapfile-ali-jetson建立8GiB、权限600并写入fstab；现有zram约3.7GiB优先级5，新磁盘swap优先级-2。总swap约11.7GiB。原fstab备份为/etc/fstab.before-ali-jetson-20261114。
\item
  软件安装与板端小模型测试使用systemd-run --user，MemoryHigh=5G、MemoryMax=6G、OOMPolicy=stop。小测试不代表大模型的加载峰值。
\item
  电源保持nvpmodel模式0（15W）。jetson\_clocks未启用；读取显示CPU为动态频率、GPU306--624.75MHz，EMC FreqOverride=0。未设置锁频或改变功耗模式。
\item
  切无桌面模式时USB网络服务停止，随后通过同一板子的局域网192.168.31.168恢复连接。板端用户csip，工作目录/home/csip/ali-jetson/deployment\_20261114。
\end{itemize}

\subsection{修复的兼容性问题}\label{s21-detail-3}

初次导入PEFT时，Transformers间接导入SciPy。Torch先加载了系统旧libstdc++，而Conda的SciPy需要GLIBCXX\_3.4.32，发生导入异常。项目专用jetson\_python.sh通过LD\_PRELOAD显式使用ali-jetson/lib/libstdc++.so.6解决；只作用于该进程，没有替换系统库或全局修改shell配置。此脚本同时限制CPU线程。

\subsection{PC预量化与跨设备核对}\label{s21-detail-4}

测试使用随机初始化的小型一层Llama测试夹具（128隐藏维度），只用于存储/内核/API兼容性，不是Qwen能力测试。

\begin{enumerate}
\def\labelenumi{\arabic{enumi}.}
\tightlist
\item
  PC在限内存systemd单元中，将FP16测试夹具量化为NF4（double\_quant=true），save\_pretrained保存。
\item
  构造非零PEFT q\_proj/v\_proj适配器，独立保存，不合并。保存同一输入的基座/适配后logits。
\item
  PC重新加载量化目录验证；传输同一量化目录、适配器与参考到Orin，直接恢复预量化权重并加载PEFT。
\end{enumerate}

\FloatBarrier\begingroup\small\setlength{\tabcolsep}{2.5pt}
\begin{longtable}[]{@{}
  >{\raggedright\arraybackslash}p{(\linewidth - 4\tabcolsep) * \real{0.4555}}
  >{\raggedleft\arraybackslash}p{(\linewidth - 4\tabcolsep) * \real{0.2722}}
  >{\raggedleft\arraybackslash}p{(\linewidth - 4\tabcolsep) * \real{0.2722}}@{}}
\caption{PC预量化与跨设备核对：设置与结果}\label{tab:s21-2}\tabularnewline
\toprule\noalign{}
\begin{minipage}[b]{\linewidth}\raggedright
指标
\end{minipage} & \begin{minipage}[b]{\linewidth}\raggedleft
PC重载
\end{minipage} & \begin{minipage}[b]{\linewidth}\raggedleft
Jetson重载
\end{minipage} \\
\midrule\noalign{}
\endfirsthead
\toprule\noalign{}
\begin{minipage}[b]{\linewidth}\raggedright
指标
\end{minipage} & \begin{minipage}[b]{\linewidth}\raggedleft
PC重载
\end{minipage} & \begin{minipage}[b]{\linewidth}\raggedleft
Jetson重载
\end{minipage} \\
\midrule\noalign{}
\endhead
\bottomrule\noalign{}
\endlastfoot
对PC保存前基座logits最大绝对差 & 0.0 & 0.0 \\
对PC保存前适配logits最大绝对差 & 0.0 & 0.0 \\
适配器引入的最大输出变化 & 0.17041015625 & 0.17041015625 \\
\end{longtable}\endgroup

PC为torch2.7.1+cu126/bnb0.48.2，Jetson为torch2.8.0/bnb0.48.0.dev0+ff389db，当前小夹具跨版本存储可恢复。零差仅对应这一固定输入的小模型，不能推断Qwen跨设备数值完全一致。实际Qwen加载、视觉前处理与图像输出精度仍需下一阶段验证。

\subsection{本地文件}\label{s21-detail-5}

\begin{itemize}
\tightlist
\item
  install\_stack.sh：固定版本安装与依赖约束。
\item
  setup\_headless\_swap.sh：用户授权后的SD swap设置。
\item
  jetson\_python.sh：进程级C++运行库修复。
\item
  stack\_probe.py：CUDA/NF4导入和矩阵测试。
\item
  smoke\_nf4.py：PC prepare / Jetson verify小夹具测试。
\item
  tiny\_nf4\_transfer/：PC夹具与PC重载结果。
\item
  software\_stack\_20261114/deployment\_20261114/：板端安装报告、pip前后清单、pip check、电源/swap状态、stack\_result.json、nf4\_transfer.log、板端verify\_result.json。
\end{itemize}


\paragraph{阶段结论}恢复后的Orin Nano 8GB建立了可运行的推理软件栈。环境与早期借用设备的测试分别记录，内存不足时通过受限进程和交换空间门控处理。 端侧可运行不等于延迟达标。使用SD卡交换空间的加载过程和无交换的稳定推理必须分开。

\stagecode{s21}

\FloatBarrier

\section{2B适配与PC、Jetson精度对照}

\label{stage:s22}

使用相同数据划分和超参数训练2B可供性LoRA，再按PC BF16、PC NF4、Jetson NF4逐步比较，区分模型规模、量化与设备的影响。

\subsection{实验设置与训练}\label{s22-detail-1}

2B基座为Qwen/Qwen3-VL-2B-Instruct，revision \path|89644892e4d85e24eaac8bacfd4f463576704203|；权重SHA256与官方文件元数据一致。直接调用既有affordance/train\_lora.py，仅替换模型源。

训练初态17--49、验证13--16的同48条样例、测试0--12的同350帧。780训练帧、600优化步、梯度累积4、学习率1e-4、r8/alpha16/dropout0.05、同任务采样权重和种子20261114，每150步验证。验证把手像素误差中位数选点；两份protocol.json逐字段相等。两者都是初态隔离、动作模板共享，不是跨场景或独立动作分布泛化。

2B选中第600步，验证误差中位数1.026像素，可训练参数1,605,632，训练记录墙钟334.7秒、峰值CUDA分配4748.9MiB。训练未读取测试集选点。

正式评测调用原eval\_qwen.py的单次main流程，沿用原适用性、提示词、解析和评分；每个完成的配置只做一遍完整测试。所有生成均为贪心、最多40新token、BF16计算和SDPA。NF4仅量化语言线性层，视觉塔和共享词嵌入的输出头保持BF16，PEFT直接加载、不合并。

4B BF16沿用已经保存的同350帧结果，并未重新挑选。NF4导出后在PC重载评测，再将同一文件和适配器送到Jetson并校验SHA256。

\subsection{精度对照}\label{s22-detail-2}

\FloatBarrier\begingroup\small\setlength{\tabcolsep}{2.5pt}
\begin{longtable}[]{@{}
  >{\raggedright\arraybackslash}p{(\linewidth - 10\tabcolsep) * \real{0.2512}}
  >{\raggedleft\arraybackslash}p{(\linewidth - 10\tabcolsep) * \real{0.1498}}
  >{\raggedleft\arraybackslash}p{(\linewidth - 10\tabcolsep) * \real{0.1498}}
  >{\raggedleft\arraybackslash}p{(\linewidth - 10\tabcolsep) * \real{0.1498}}
  >{\raggedleft\arraybackslash}p{(\linewidth - 10\tabcolsep) * \real{0.1498}}
  >{\raggedleft\arraybackslash}p{(\linewidth - 10\tabcolsep) * \real{0.1498}}@{}}
\caption{精度对照：设置与结果}\label{tab:s22-1}\tabularnewline
\toprule\noalign{}
\begin{minipage}[b]{\linewidth}\raggedright
配置（均512像素）
\end{minipage} & \begin{minipage}[b]{\linewidth}\raggedleft
把手严格命中
\end{minipage} & \begin{minipage}[b]{\linewidth}\raggedleft
4像素邻域命中
\end{minipage} & \begin{minipage}[b]{\linewidth}\raggedleft
把手3D中位误差/cm
\end{minipage} & \begin{minipage}[b]{\linewidth}\raggedleft
柜体框IoU
\end{minipage} & \begin{minipage}[b]{\linewidth}\raggedleft
有效3D样本
\end{minipage} \\
\midrule\noalign{}
\endfirsthead
\toprule\noalign{}
\begin{minipage}[b]{\linewidth}\raggedright
配置（均512像素）
\end{minipage} & \begin{minipage}[b]{\linewidth}\raggedleft
把手严格命中
\end{minipage} & \begin{minipage}[b]{\linewidth}\raggedleft
4像素邻域命中
\end{minipage} & \begin{minipage}[b]{\linewidth}\raggedleft
把手3D中位误差/cm
\end{minipage} & \begin{minipage}[b]{\linewidth}\raggedleft
柜体框IoU
\end{minipage} & \begin{minipage}[b]{\linewidth}\raggedleft
有效3D样本
\end{minipage} \\
\midrule\noalign{}
\endhead
\bottomrule\noalign{}
\endlastfoot
工作站 / 2B / 基座 / bf16 & 0.0\% & 0.0\% & --- & 0.000 & 0/350 \\
工作站 / 2B / 低秩适配 / bf16 & 93.7\% & 98.6\% & 1.515 & 0.987 & 350/350 \\
工作站 / 4B / 基座 / bf16 & 5.7\% & 17.7\% & 7.798 & 0.647 & 350/350 \\
工作站 / 4B / 低秩适配 / bf16 & 92.3\% & 98.6\% & 1.368 & 0.986 & 350/350 \\
工作站 / 2B / 低秩适配 / nf4 & 91.1\% & 95.7\% & 1.551 & 0.976 & 350/350 \\
工作站 / 4B / 低秩适配 / nf4 & 92.3\% & 99.1\% & 1.340 & 0.981 & 350/350 \\
Jetson / 2B / 低秩适配 / nf4 & 91.1\% & 95.7\% & 1.548 & 0.975 & 350/350 \\
\end{longtable}\endgroup

\stagefigure{figures/ch07/s22_precision_comparison.png}{可供性定位精度对照}{fig:s22-detail-1}

严格命中要求预测像素本身落入目标掩码；邻域命中为半径4像素的方形窗口与掩码相交。解析失败按失败计数，不从命中率分母剔除；表中3D只在可投影输出上计算，同时列出覆盖数。

3D来自仿真真值深度与相机参数反投影，仍有约1.4cm表面点/几何中心误差下限，不能解释为真实测距或抓取精度。完整碗/放置点指标见SUMMARY.json；350帧包含同回合相关画面，不能当作350次独立机器人任务。

\subsection{输出协议与诊断}\label{s22-detail-3}

2B基座在长定位提示词下常输出额外label说明，在40-token上限内截断：把手350条、放置点143条、柜体框350条均解析失败，碗171条能正常解析。因此基座正式表中零命中是当前接口不合格，不能直接解释为完全没有视觉定位能力。

为了区分格式与定位，另在CPU读取保存的原始文本，提取其中已经完整输出的坐标数组，没有重调模型。这个非正式诊断得到：基座把手严格命中1.7\%、邻域命中6.3\%、3D中位误差16.11cm，柜体IoU0.594。诊断不替代生产解析，也没有参与检查点选择。LoRA收益同时包含输出格式对齐和定位改善，不能将全部收益归于视觉能力。

\subsection{PC/Jetson逐条一致性}\label{s22-detail-4}

2B NF4对照共1014条：原始文本完全相同774条（76.33\%），解析后的坐标/框完全相同774条（76.33\%）。实际输入token、图像张量和网格哈希的匹配数量为输入词元、图像张量和图像网格各1014条完全一致，输出词元774条一致。

相同输入不意味着不同GPU内核必须逐bit相同；若有不同，保留逐条输出。PC与板端Transformers/PEFT版本相同，Torch、torchvision和bitsandbytes使用各自兼容构建，故平台对照也包含软件栈差异，不能把所有差异只归因于GPU硬件。量化前后与设备变化分步比较，不混为一个变化。该阶段四个PC配置的1014条实际输入token、图像张量和网格均逐条相同，记录在PC\_INPUT\_CHECK.json；历史4B BF16未保存同样的逐条输入元数据，只能核对原始协议、图片与提示词，不能补造历史哈希。

\subsection{加载与内存}\label{s22-detail-5}

2B板端该次进程首次加载（处理器、模型、适配器及GPU同步）17.21秒。操作系统文件缓存未清空，不是冷盘启动基准。

进程VmHWM峰值1.625GiB，稳定期cgroup采样最高1.854GiB，CUDA分配峰值2.111GiB，CUDA保留峰值2.168GiB；tegrastats全阶段整机RAM最高4089MB。以上统计口径不能相加；整机RAM包含系统和缓存，CUDA内存也可能与共享物理内存重叠。

板端内核不提供cgroup memory.peak，故该字段为null，没有拿稳定期采样最大值冒充完整加载峰值。保留VmHWM、CUDA峰值和从加载前开始的tegrastats原始记录。

正式评测模型cgroup最大swap为0字节，稳定段系统交换I/O页变化为换入和换出页数增量均为0；结束时/proc/swaps没有活动设备。加载时保留交换空间，正式评测前swapoff，结束后恢复原zram及8GiB SD swap。安装与推理均有systemd内存上限，板端MemoryMax=6G，15W且未锁jetson\_clocks。

4B板端全量350帧未执行；本阶段的完整板端精度与逐条一致率结论仅属于2B。 可选加载探针已通过，加载39.52秒，并在验证集图片上完成一次查询；这不是4B准确率或延迟评测。

\subsection{工程问题与处理}\label{s22-detail-6}

PC第一次NF4导出出现CPU内存膨胀并触发systemd保护，训练和两份BF16完整评测未丢失、未重跑。之后改用独立受限单元，并显式保留共享输出头，成功保存NF4文件；后续限制编译线程并逐配置运行，不提高内存上限。PC驱动脚本因运行中编辑发生收尾解析错误，但四份该阶段PC评测均已保存1014条唯一问答和完整资源记录；驱动错误单列留档，没有重跑测试或补造输出。

板端第一次启动推理虽然模型自身swap为0，系统仍有交换I/O。31条未完成尝试标为无效。正式新tag增加加载/无swap门控，swapoff确认后才开始350帧单遍评测。原始失败文件仍保留，不以成功结果覆盖。

保存后的tokenizer曾给出通用正则告警。该阶段没有根据告警擅自改tokenizer；实际输入token和图像哈希用于核实对照输入。正式结论只限当前图像、提示词、单场景和检查点。


\paragraph{阶段结论}2B BF16严格把手命中328/350；2B NF4在PC与Jetson上均为319/350。全部1014次问答的输入哈希一致，原文和坐标输出一致774次。 BF16是脑浮点16位格式，NF4是4位正规化量化格式。稳定评测无交换空间活动，设备输出一致率不等于总体精度。

\stagecode{s22}

\FloatBarrier

\section{Jetson视觉服务与PC仿真联调}

\label{stage:s23}

将2B量化模型常驻板端，通过SSH隧道向PC提供把手点。PC保持唯一物理场景，完成导航、反投影、规划和关节执行。

\subsection{视觉服务与PC联合仿真}\label{s23-detail-1}

服务只绑定Jetson的127.0.0.1，通过SSH隧道向PC的127.0.0.1转发，串行处理请求，不在PC加载Qwen。使用512尺寸，因为该配置已有上一阶段350帧精度核对；未选择尚未验证精度的新尺寸。服务返回二维把手点，PC用仿真深度反投影并交给MoveIt实际驱动同一世界中的关节。

机器人在模型返回期间暂停仿真推进，所以联合回合通过只能说明功能链路可用，不能替代\(\leq\)3秒实时要求。联合回合只做预接触，不执行抓取或原LIBERO整任务。

当前记录3回合，预接触通过3/3。

\FloatBarrier\begingroup\small\setlength{\tabcolsep}{2.5pt}
\begin{longtable}[]{@{}
  >{\raggedleft\arraybackslash}p{(\linewidth - 8\tabcolsep) * \real{0.2000}}
  >{\raggedright\arraybackslash}p{(\linewidth - 8\tabcolsep) * \real{0.2000}}
  >{\raggedright\arraybackslash}p{(\linewidth - 8\tabcolsep) * \real{0.2000}}
  >{\raggedleft\arraybackslash}p{(\linewidth - 8\tabcolsep) * \real{0.2000}}
  >{\raggedleft\arraybackslash}p{(\linewidth - 8\tabcolsep) * \real{0.2000}}@{}}
\caption{视觉服务与PC联合仿真：设置与结果}\label{tab:s23-1}\tabularnewline
\toprule\noalign{}
\begin{minipage}[b]{\linewidth}\raggedleft
初态
\end{minipage} & \begin{minipage}[b]{\linewidth}\raggedright
导航成功
\end{minipage} & \begin{minipage}[b]{\linewidth}\raggedright
预接触通过
\end{minipage} & \begin{minipage}[b]{\linewidth}\raggedleft
真值误差 mm
\end{minipage} & \begin{minipage}[b]{\linewidth}\raggedleft
视觉残差 mm
\end{minipage} \\
\midrule\noalign{}
\endfirsthead
\toprule\noalign{}
\begin{minipage}[b]{\linewidth}\raggedleft
初态
\end{minipage} & \begin{minipage}[b]{\linewidth}\raggedright
导航成功
\end{minipage} & \begin{minipage}[b]{\linewidth}\raggedright
预接触通过
\end{minipage} & \begin{minipage}[b]{\linewidth}\raggedleft
真值误差 mm
\end{minipage} & \begin{minipage}[b]{\linewidth}\raggedleft
视觉残差 mm
\end{minipage} \\
\midrule\noalign{}
\endhead
\bottomrule\noalign{}
\endlastfoot
2 & 是 & 是 & 18.2 & 5.3 \\
3 & 是 & 是 & 8.5 & 5.8 \\
4 & 是 & 是 & 15.4 & 7.5 \\
\end{longtable}\endgroup

服务稳定段无swap检查：True。模型首加载11.77秒。


\paragraph{阶段结论}三个联合预接触回合3/3通过，误差18.2、8.5、15.4毫米；六次正式响应的图片哈希与PC源帧匹配，稳定段无交换活动。 模型返回期间仿真暂停，联合成功不能替代实时达标。512像素配置沿用既有350帧精度对照，未做实体机械臂联调。

\stagecode{s23}

\stagevideo{videos/ch07/jetson_vision_pc_execution.mp4}

\FloatBarrier

\section{端侧时延、资源与算子剖析}

\label{stage:s24}

固定验证集五张图，在15瓦功耗下比较512、384和256像素，记录输出长度、墙钟时延和资源；另做一次独立的算子性能剖析。

\subsection{时延设置与结果}\label{s24-detail-1}

输入为冻结验证集排序前5帧，未使用test350。每帧询问handle、bowl、place、cabinet四类问题，每档20次，另各暖机一次不计入。尺寸按图像轮换顺序，减少固定顺序影响；P95只是探索性小样本分位数。max\_new\_tokens为40，贪心解码，未缩短原提示词。

墙钟时间从PC/板端测试脚本读取和处理图像开始，到生成文本结束；不含网络传输。模型内时间另存，不能将其替代端到端时延。输出长度一并报告，因为不同问题生成token数不同。

\FloatBarrier\begingroup\small\setlength{\tabcolsep}{2.5pt}
\begin{longtable}[]{@{}
  >{\raggedleft\arraybackslash}p{(\linewidth - 12\tabcolsep) * \real{0.1429}}
  >{\raggedright\arraybackslash}p{(\linewidth - 12\tabcolsep) * \real{0.1429}}
  >{\raggedleft\arraybackslash}p{(\linewidth - 12\tabcolsep) * \real{0.1429}}
  >{\raggedleft\arraybackslash}p{(\linewidth - 12\tabcolsep) * \real{0.1429}}
  >{\raggedleft\arraybackslash}p{(\linewidth - 12\tabcolsep) * \real{0.1429}}
  >{\raggedleft\arraybackslash}p{(\linewidth - 12\tabcolsep) * \real{0.1429}}
  >{\raggedright\arraybackslash}p{(\linewidth - 12\tabcolsep) * \real{0.1429}}@{}}
\caption{时延设置与结果：设置与结果}\label{tab:s24-1}\tabularnewline
\toprule\noalign{}
\begin{minipage}[b]{\linewidth}\raggedleft
尺寸
\end{minipage} & \begin{minipage}[b]{\linewidth}\raggedright
查询
\end{minipage} & \begin{minipage}[b]{\linewidth}\raggedleft
次数
\end{minipage} & \begin{minipage}[b]{\linewidth}\raggedleft
P50 秒
\end{minipage} & \begin{minipage}[b]{\linewidth}\raggedleft
P95 秒
\end{minipage} & \begin{minipage}[b]{\linewidth}\raggedleft
平均输出token
\end{minipage} & \begin{minipage}[b]{\linewidth}\raggedright
\(\leq\)3秒P95
\end{minipage} \\
\midrule\noalign{}
\endfirsthead
\toprule\noalign{}
\begin{minipage}[b]{\linewidth}\raggedleft
尺寸
\end{minipage} & \begin{minipage}[b]{\linewidth}\raggedright
查询
\end{minipage} & \begin{minipage}[b]{\linewidth}\raggedleft
次数
\end{minipage} & \begin{minipage}[b]{\linewidth}\raggedleft
P50 秒
\end{minipage} & \begin{minipage}[b]{\linewidth}\raggedleft
P95 秒
\end{minipage} & \begin{minipage}[b]{\linewidth}\raggedleft
平均输出token
\end{minipage} & \begin{minipage}[b]{\linewidth}\raggedright
\(\leq\)3秒P95
\end{minipage} \\
\midrule\noalign{}
\endhead
\bottomrule\noalign{}
\endlastfoot
512 & all & 20 & 5.06 & 7.15 & 19 & 未达标 \\
512 & handle & 5 & 5.07 & 5.09 & 17 & 未达标 \\
512 & cabinet & 5 & 7.14 & 7.15 & 25 & 未达标 \\
384 & all & 20 & 4.91 & 6.97 & 19 & 未达标 \\
384 & handle & 5 & 4.92 & 4.95 & 17 & 未达标 \\
384 & cabinet & 5 & 6.97 & 6.99 & 25 & 未达标 \\
256 & all & 20 & 4.78 & 6.85 & 19 & 未达标 \\
256 & handle & 5 & 4.79 & 4.81 & 17 & 未达标 \\
256 & cabinet & 5 & 6.83 & 6.86 & 25 & 未达标 \\
\end{longtable}\endgroup

首次加载24.34秒，未清除OS页缓存，不称冷盘启动。稳定段模型swap与全局交换I/O均为0。NF4保持视觉塔及共享lm\_head为BF16，其余量化设置与上一阶段一致。

\subsection{功耗、频率与资源}\label{s24-detail-2}

使用15W模式0，未调用jetson\_clocks锁频。实测GPU下限306MHz、上限624.75MHz，nvhost\_podgov动态调频；CPU0上下限729.6/1510.4MHz。tegrastats逐100ms记录实际负载与频率。本板当前启动配置不提供MAXN SUPER模式，该次没有切换或刷写固件。

模型进程MemoryHigh5GiB、MemoryMax6GiB。加载时允许SD/zram作储备，加载完成后关闭swap再开始稳定测量；结束恢复原SD和zram。不能把SD交换造成的慢速结果算作有效实时测试。

缩图后时延仍接近5秒，且25-token的柜体框输出明显慢于17-token的点输出，说明继续只减像素的收益有限。时延与输出长度相关只是观察结果，不能代替CUDA profiler对预填充、逐token解码和内核开销的分解；不据此武断归因。

\subsection{单次算子剖析}\label{s24-detail-3}

在正式时延之外，另用一张验证图做诊断，生成17个token，捕获41434个CUDA内核事件。内核时长累加2.85秒；此值可能含重叠，不能与CPU时间直接相加。开启profiler后的推理时间不混入上表P50/P95。

普通用户首次受CUPTI权限限制，没有GPU事件；保留该记录。随后只给这一受限剖析进程管理员权限，没有修改驱动全局权限。第一次管理员启动因门控文件权限未开始剖析，另存失败；修复门控后完成。

按CPU表，cudaLaunchKernel共41431次、自耗时约1.08秒；4bit矩阵向量计算约1.11秒CUDA时间，另有解量化、拷贝与大量小算子。量化解码与调度开销是下一步候选优化点，但不能直接扣掉这些时间宣称已达标。当前ali-jetson未安装Triton。已找到官方Jetson AI Lab匹配的3.4.0预编译包，但安装前核对到当前Transformers的NF4量化器没有覆盖is\_compileable=False，现有自动编译入口会拒绝该组合；因此没有安装或绕过检查。该限制不等于所有NF4模型都无法手工编译，需另做兼容性验证。

后续可在独立环境验证静态KV缓存与编译解码，先确认输出一致，再测速度和完整精度；不再仅靠换更大模型或缩图。接口参考\href{https://docs.pytorch.org/docs/2.8/profiler.html}{PyTorch 2.8 profiler}、\href{https://developer.nvidia.com/nvidia-development-tools-solutions-err-nvgpuctrperm-cupti}{NVIDIA CUPTI权限说明}、\href{https://huggingface.co/docs/transformers/v4.57.1/en/llm_optims}{Transformers静态缓存与编译}。


\subsection{资源统计的不同口径}\label{s24-detail-4}

资源摘要（MiB，不同口径不能相加，Jetson使用共享物理内存）：

\FloatBarrier\begingroup\small\setlength{\tabcolsep}{2.5pt}
\begin{longtable}[]{@{}
  >{\raggedright\arraybackslash}p{(\linewidth - 8\tabcolsep) * \real{0.1946}}
  >{\raggedleft\arraybackslash}p{(\linewidth - 8\tabcolsep) * \real{0.1946}}
  >{\raggedleft\arraybackslash}p{(\linewidth - 8\tabcolsep) * \real{0.1946}}
  >{\raggedleft\arraybackslash}p{(\linewidth - 8\tabcolsep) * \real{0.1946}}
  >{\raggedleft\arraybackslash}p{(\linewidth - 8\tabcolsep) * \real{0.2215}}@{}}
\caption{资源统计的不同口径：设置与结果}\label{tab:s24-2}\tabularnewline
\toprule\noalign{}
\begin{minipage}[b]{\linewidth}\raggedright
阶段
\end{minipage} & \begin{minipage}[b]{\linewidth}\raggedleft
进程VmHWM
\end{minipage} & \begin{minipage}[b]{\linewidth}\raggedleft
cgroup采样最大
\end{minipage} & \begin{minipage}[b]{\linewidth}\raggedleft
CUDA分配峰值
\end{minipage} & \begin{minipage}[b]{\linewidth}\raggedleft
tegrastats整板RAM最大
\end{minipage} \\
\midrule\noalign{}
\endfirsthead
\toprule\noalign{}
\begin{minipage}[b]{\linewidth}\raggedright
阶段
\end{minipage} & \begin{minipage}[b]{\linewidth}\raggedleft
进程VmHWM
\end{minipage} & \begin{minipage}[b]{\linewidth}\raggedleft
cgroup采样最大
\end{minipage} & \begin{minipage}[b]{\linewidth}\raggedleft
CUDA分配峰值
\end{minipage} & \begin{minipage}[b]{\linewidth}\raggedleft
tegrastats整板RAM最大
\end{minipage} \\
\midrule\noalign{}
\endhead
\bottomrule\noalign{}
\endlastfoot
时延筛查 & 1657 & 2310 & 2162 & 4060 \\
服务联调 & 1660 & 1661 & 2162 & 4063 \\
\end{longtable}\endgroup


\paragraph{阶段结论}三档P50约5.06、4.91、4.78秒，P95约7.15、6.97、6.85秒，均未达3秒目标。单次剖析捕获约4.1万个CUDA内核事件。 P50和P95分别为第50、95百分位数，样本量小。剖析开销不混入时延；现有量化器自动编译兼容性尚未打通，没有为新尺寸追加全量精度评测。

\stagecode{s24}

\FloatBarrier
